\chapter{Fazit und Ausblick}
\label{ch:fazit}

\section{Fazit}
\label{sec:fazit}

Durch die Umstellung des Kühlkreislauflayouts konnte erreicht werden, dass
alle vier Motoren annähernd dieselbe Eintritts- und Austrittstemperatur des
Kühlwassers aufweisen. Dadurch werden die Motoren gleichmäßiger gekühlt. Im
Endurance traten seitdem keine Motorausfälle durch Überhitzung mehr auf
\cite{TODO_Quelle_bzw_Verweis_auf_Messdaten}. Das neue Entlüftungskonzept
verkürzt außerdem die Zeit für das Befüllen und Entlüften des Kühlkreislaufs
während der Events.

% TODO: Satz(e) zur Gesamtbilanz ergänzen: Was leistet das Simulink-Modell
% (Kapitel~\ref{ch:systemmodelle}), wie weit ist die Sensorik validiert
% (Abschnitt~\ref{subsec:sensorik}), was ist der Stand des Testplans
% (Abschnitt~\ref{sec:Testplanung})? Der Fazit-Text bezieht sich bisher nur auf
% die Layout-Änderung.

\section{Ausblick}
\label{sec:ausblick}

Für die kommende Saison sollte ein U-Radiator in Betracht gezogen werden.
Dieser verursacht zwar einen höheren Druckverlust, kann das Kühlmittel
jedoch stärker abkühlen und damit die Kühlleistung des Gesamtsystems
steigern. Der höhere Druckverlust ist bei der Pumpenauslegung zu
berücksichtigen und kann mit dem entwickelten Systemmodell vorab untersucht
werden.

Langfristig soll eine effizientere Kombination aus Motor und Inverter im
Fahrzeug eingesetzt werden. Beispielsweise erzeugen die Inverter des
Teams aus Dresden~\cite{TODO_Quelle} den sinusförmigen Stromverlauf mit einer
deutlich höheren Taktfrequenz. Das reduziert die Stromwelligkeit und damit
die Verluste, sodass diese Inverter besonders im Hochdrehzahlbereich
effizienter arbeiten. Mit einer solchen Kombination sinkt die abzuführende
Verlustleistung. Dadurch wären dünnere Schläuche und kleinere Pumpen
möglich, was das Gewicht des Kühlsystems deutlich verringern würde.

\pagebreak